\chapter{Conclusion and Future Work} \label{chapter:ch7_conclusion}

\section{Conclusion}

This thesis asked whether collaborative-robot tasks that change frequently, represented here by pick-and-place and seam welding, can be authored more directly by specifying spatial intent in the real workspace instead of relying only on teach pendant, offline programming, or hand-guiding workflows. Within prototype scope, the answer is positive for selected task-oriented workflows. The developed system shows that poses and trajectories entered in the workspace can be combined with scene detection, task-specific interpretation, and robot execution to produce executable procedures on a real robot without exposing the operator to a full robot programming environment.

The main contribution is a reusable task-oriented authoring path rather than only a tracked input device or two isolated demonstrations. Within a selected use case, the same path supports authoring, confirmation into executable commands, execution, visualization, and persistence. A shared runtime with explicit interfaces allows both pick-and-place and seam welding to reuse the same infrastructure.

The result nevertheless remains a proof of concept. Most development and testing took place in the REMII offline environment, physical validation is limited, and the thesis does not include a formal user study, a baseline comparison against teach pendant or offline programming workflows, or a quantitative characterization of tracking accuracy, repeatability, or productivity. Within that boundary, seam welding is the clearest practical fit because authored trajectory input maps directly to a welding path. Pick-and-place mainly shows that an authored pose can be combined with refreshed scene information at execution time, but it depends more strongly on the narrow scene model used in the prototype.

\section{Near-Term Future Work}

Improving sensing, tracking, and calibration is the most direct near-term step. More robust pen tracking, better calibration diagnostics, fuller use of pen-side IMU data, and a richer scene estimate than the current fiducial-box representation would directly improve the present workflow. The same work should also support bounded characterization of tracking stability, scene uncertainty, and repeated scene acquisition in the real workspace.

Execution support also remains incomplete at the current prototype level. A clearer abstraction for process tools, explicit motion validation and collision checking for authored tasks, and simulated execution for safer inspection before real-robot use would make the current workflow easier to trust. For tasks that do not depend on runtime scene adaptation, it is also worth exploring a path from authored commands to a lower-level robot-side export or compilation target.

Separate from those engineering steps, the prototype still needs broader validation and limited platform completion. Broader physical testing and bounded characterization of repeatability, tracking accuracy, and scene-estimation stability would make the present evidence easier to interpret. Some remaining cross-platform work and compact interface cleanup belong here as secondary tasks, because they affect whether the current prototype can be used reliably in repeated experiments.

\section{Longer-Term Research Directions}

Because the architecture is not tied to the current pen-based interface, richer input modalities remain a natural research direction. Hand tracking, VR or AR controllers, and combinations of speech and gesture could be explored within the same task-oriented model. This would help clarify which forms of spatial input remain natural for different task families without losing control over the resulting robot behavior.

A separate question is how this authoring model should connect to existing robot-programming practice. Spatial authoring could be integrated more directly with teach-pendant or controller workflows instead of remaining a separate authoring tool, and that integration should be studied together with a formal user study. Those steps would help determine when the approach is actually useful, where its limits remain too restrictive, and how operators understand the boundary between selected use cases and more general robot programming.

Only after that should the prototype be pushed across a broader range of task families and operating conditions. Richer welding scenarios, additional fabrication or surface-processing tasks, and more varied human-in-the-loop workflows would provide a stronger basis for judging that applicability boundary. The present prototype is a bounded starting point for that work, but broader claims still depend on technical completion and evaluation on concrete task families.
